% Source-derived chapter 07: The asymptotic Putman-Wieland conjecture
% Original source: canonical-representations-surface-groups
\chapter{The asymptotic Putman-Wieland conjecture}
\label{canonical-representations-surface-groups-chapter-07}
\source{canonical-representations-surface-groups}

\begin{remark}[Source section]
\label{canonical-representations-surface-groups-chapter-07-source}
\source{canonical-representations-surface-groups}
\source{canonical-representations-surface-groups}
This chapter reproduces the corresponding section of the retrieved
LaTeX source, including its definitions, statements, examples, and
proofs. It has not yet been translated into Lean.
\notready
\end{remark}

% The source section title was: The asymptotic Putman-Wieland conjecture
\label{section:Putman-Wieland}
\source{canonical-representations-surface-groups}
\subsection{The statement of the Putman-Wieland Conjecture}
\label{subsection:geometric-setup}
\source{canonical-representations-surface-groups}


We next discuss some applications of our methods to the Putman-Wieland
conjecture (\autoref{conjecture:putman-wieland}), a major open problem in geometric topology. Much of the interest in the Putman-Wieland conjecture arises from its close relationship to Ivanov's conjecture that mapping class groups $\on{Mod}_{g,n}$ do not virtually surject onto $\mathbb{Z}$ for $g\gg 0$ \cite[\S7]{Ivanov:problems}, as explained in
\cite[Theorem 1.3]{putmanW:abelian-quotients}. In particular, the Putman-Wieland conjecture for $g\gg 0$ is equivalent to Ivanov's conjecture for $g\gg 0$.
\begin{notation}
\source{canonical-representations-surface-groups}
\notready
	\label{notation:prym}
\source{canonical-representations-surface-groups}
Let $\Sigma_g$ be an oriented surface of genus $g \geq 0$, and let
$\Sigma_{g,b,p}$ be the complement of $b\geq 0$ disjoint open discs and $p\geq 0$ disjoint points in $\Sigma_g$, so that $\Sigma_{g,b,p}$ is an oriented genus $g$ surface with $b$ boundary components and $p$ punctures. We let $\on{PMod}_{g,b,p}$ be the subgroup of the mapping class group of $\Sigma_{g,b,p}$ fixing the punctures and boundary components pointwise.
%Recall a subgroup $G' \subset G$ is characteristic if $f(G') \subset G'$ for all
%$f \in \on{Aut}(G)$.
Fix a basepoint $v_0 \in \Sigma_{g,b,p}$, which we count as an additional puncture to
obtain an action of $\on{PMod}_{g,b,p+1}$ on $\pi_1(\Sigma_{g,b,p},v_0)$.
Let $K \trianglelefteq \pi_1(\Sigma_{g,b,p}, v_0)$ be a finite index normal subgroup
corresponding to a finite Galois covering space $\Sigma_{g',b',p'} \to \Sigma_{g,b,p}$,
or equivalently a surjection
$\phi: \pi_1(\Sigma_{g,b,p}, v_0) \twoheadrightarrow H$, with $H=\pi_1(\Sigma_{g,b,p}, v_0)/K$ a finite group.
Let $\Gamma \subset \on{PMod}_{g,b,p+1}$ denote the finite index subgroup preserving
$\phi$ up to conjugacy.
Viewing
$H_1(\Sigma_{g',b',p'},\mathbb C)$ as an $H$-representation,
if $\rho$ is an irreducible $H$-representation we let 
\begin{align*}
	H_1(\Sigma_{g',b',p'},\mathbb C)^\rho := \rho \otimes \Hom_H(\rho,
H_1(\Sigma_{g',b',p'},\mathbb C))
\end{align*}
denote the $\rho$-isotypic component.
%, i.e.~the span of all $H$-subrepresentations of $H_1(\Sigma_{g',b',p'},\mathbb C)$ isomorphic to $\rho$.
We obtain an action of $\Gamma$ on $K^{\text{ab}}\otimes \mathbb{C}=H_1(\Sigma_{g',b',p'},\mathbb C)$
and hence on the characteristic subrepresentation
$H^1(\Sigma_{g',b',p'},\mathbb C)^\rho$.
By filling in the punctures and deleted discs, we also obtain
an action of $\Gamma$ on $V_K := H_1(\Sigma_{g'}, \mathbb C)$, referred to in 
\cite[p. 80-81]{putmanW:abelian-quotients} as a {\em higher Prym
representation}.
\end{notation}

\begin{definition}
\source{canonical-representations-surface-groups}
\notready
	\label{definition:pw}
\source{canonical-representations-surface-groups}
	Fix a finite group $H$ and nonnegative integers $g, b,$ and $p$.
	Let $\PW_{g,b,p}^H$ be the statement that for any
	surjection $\phi: \pi_1(\Sigma_{g,b,p}, v_0) \twoheadrightarrow
	H$, taking $K := \ker \phi$ and $v \in V_K$ any nonzero vector,
	$v$ has infinite orbit under $\Gamma$.
\end{definition}

\begin{conjecture}
\source{canonical-representations-surface-groups}
\notready[Putman-Wieland, ~\protect{\cite[Conjecture
	1.2]{putmanW:abelian-quotients}}]
	\label{conjecture:putman-wieland}
\source{canonical-representations-surface-groups}
	$\PW_{g,b,p}^H$ holds for every group $H$ with $g \geq 2, b \geq 0, p \geq 0$.
\end{conjecture}

\begin{remark}
\source{canonical-representations-surface-groups}
\notready
\label{remark:putman-wieland-differences}
\source{canonical-representations-surface-groups}
We note that 
\cite[Conjecture 1.2]{putmanW:abelian-quotients}
is stated without a group $H$ and with $\mathbb Q$ coefficients instead of
$\mathbb C$ coefficients.
We use this equivalent statement to more easily state our results,
which imply that 
$\PW_{g,b,p}^H$ holds whenever $\# H$ is small compared to the genus $g$. 
\end{remark}
%Putman-Wieland also analyze cohomology with rational coefficients, but their conjecture remains unchanged after extending scalars to $\mathbb{C}$. We do so, as it will be convenient for all irreducible representations to be absolutely irreducible.
\begin{remark}
\source{canonical-representations-surface-groups}
\notready
	\label{remark:}
\source{canonical-representations-surface-groups}
	There is a counterexample to the Putman-Wieland conjecture when $g = 2$ \cite[Theorem 1.3]{markovic}.
\end{remark}


\subsection{Our results}

We are able to verify $\PW_{g,b,p}^H$ in many new cases.
Moreover, we prove the following more general statement, ruling out
not only invariant vectors in $H^1(\Sigma_{g'}, \mathbb C)$, but even ruling out low-dimensional invariant subspaces
of isotypic components.
The proof is given below in \autoref{subsection:proof-putman-wieland}.
\begin{theorem}
\source{canonical-representations-surface-groups}
\notready
	\label{theorem:asymptotic-putman-wieland}
\source{canonical-representations-surface-groups}
	With notation as in \autoref{notation:prym}, let $\rho$ be an irreducible
	complex $H$-representation and $\Gamma' \subset \Gamma$ be a finite-index subgroup. Then
	$H^1(\Sigma_{g',b',p'},\mathbb C)^\rho$
	has no nonzero $\Gamma'$-invariant subrepresentations 
	of dimension strictly less than $2g - 2\dim \rho$.
	The same holds for
	$H_1(\Sigma_{g'},\mathbb C)^\rho$ in place of 
	$H^1(\Sigma_{g',b',p'},\mathbb C)^\rho$.
\end{theorem}

\begin{remark}
\source{canonical-representations-surface-groups}
\notready
In \autoref{theorem:asymptotic-putman-wieland}, it is important that we restrict our attention to subrepresentations, as opposed to arbitrary subquotients, as the $\Gamma$-representation 	$H^1(\Sigma_{g',b',p'},\mathbb C)$ is not in general semisimple for $b'+p'>0$.
\end{remark}

A perhaps more concrete corollary is the following.
\begin{corollary}
\source{canonical-representations-surface-groups}
\notready
	\label{corollary:rep-theory-PW}
\source{canonical-representations-surface-groups}
	$\PW_{g,b,p}^H$ holds for every group $H$ 
	such that every irreducible representation $\rho$ of $H$ has dimension
	$\dim \rho < g$.
\end{corollary}
\begin{proof}
	Fix an irreducible representation $\rho$ of $H$. By assumption, $\dim \rho < g$. Then
	$2g - 2\dim \rho > 1$, and so by
	\autoref{theorem:asymptotic-putman-wieland},
	$H_1(\Sigma_{g'},\mathbb C)^\rho$ has no nonzero $\Gamma'$-invariant $1$-dimensional subrepresentations for any $\rho$ and any finite index $\Gamma'\subset\Gamma$.
	In particular, $H_1(\Sigma_{g'},\mathbb C)$ has no nonzero invariant
	vectors under any finite index $\Gamma'$. But if there was a nonzero vector $v\in V_K$ with finite orbit, its stabilizer would yield a finite index $\Gamma'$ with a fixed vector.
\end{proof}
Even more concretely, we can give the following bound on $\#H$ independent of its
representation theory.
\begin{corollary}
\source{canonical-representations-surface-groups}
\notready
	\label{corollary:asymptotic-putman-wieland}
\source{canonical-representations-surface-groups}
	For fixed $g, b, p \geq 0$,
	$\PW_{g,b,p}^H$ holds for every group $H$ with $\# H < g^2$.
\end{corollary}
\begin{proof}
	If $H$ is a finite group of order $\# H < g^2$, then because $\#
	H$ is the sum of the squares of the dimensions of the irreducible representations of
	$H$ by \cite[(2.19)]{fultonH:representation-theory},
	no irreducible representations of $H$ have
	dimension $\geq g$. The result follows
	from \autoref{corollary:rep-theory-PW}.
\end{proof}

In order to prove \autoref{theorem:asymptotic-putman-wieland},
we first note that it suffices to consider the case of surfaces without boundary.
\begin{lemma}
\source{canonical-representations-surface-groups}
\notready
	\label{lemma:ignore-boundary}
\source{canonical-representations-surface-groups}
	With notation as in \autoref{notation:prym},
	for any $g \geq 0, b, p \geq 0$, 
	and $\rho$ an irreducible $H$-representation, let $\Gamma'\subset
	\Gamma$ be a finite index subgroup. Then there exists a finite index
	$\Gamma''\subset \on{PMod}_{g,0,b+p+1}$ and 
 an isomorphism	
$H_1(\Sigma_{g',b',p'}, \mathbb C)^\rho \to H_1(\Sigma_{g,0,b'+p'}, \mathbb
C)^\rho$
inducing a bijection from $\Gamma'$-invariant subspaces to
$\Gamma''$-invariant subspaces.
\end{lemma}
\begin{proof}
% Let $\Gamma\subset \pi_1(\Sigma_{g,b,p})=\pi_1(\Sigma_{g, 0, b+p})$ be a
% characteristic subgroup of finite index, corresponding to some ramified cover
% $\Sigma_{g'}\to \Sigma_g$. We wish to show that the $\on{Mod}_{g, 0,
% b+p}$-action on $H_1(\Sigma_{g'}, \mathbb{Q})$ admits no nonzero vectors with
% finite orbit if and only if the same is true of the action of $\on{Mod}_{g, b, p}$. 
%
%Let $\Sigma_{g',b',p'} \subset \Sigma_{g'}$ denote the preimage of
%$\Sigma_{g,b,p} \subset \Sigma_{g}$ and $\Sigma_{g',b'+p'} \subset \Sigma_{g'}$ 
%the preimage of $\Sigma_{g,0,b+p} \subset \Sigma_g$.
Define the natural inclusion $\iota: \Sigma_{g',b',p'} \hookrightarrow
\Sigma_{g,0,b'+p'}$,
shrinking boundary components to punctures.
There is a deformation retraction from $\Sigma_{g,0,b'+p'}$ onto the subspace
$\iota(\Sigma_{g',b',p'})$
and so $\iota$ induces an isomorphism
$H_1(\Sigma_{g',b',p'}, \mathbb C) \to H_1(\Sigma_{g,0,b'+p'}, \mathbb C).$
The kernel of the induced surjection $\on{PMod}_{g, b,p} \to \on{PMod}_{g,0,b+p}$
is generated by $\mathbb Z^b$ acting on $S_{g,b,p}$ by rotating the
boundary components,
\cite[Theorem 3.18]{farbM:a-primer}
and hence this $\mathbb Z^b$ acts trivially on 
$H_1(\Sigma_{g',b',p'}, \mathbb C)$. 
Let $\Gamma''$ be the image of $\Gamma'$ under this natural map $\on{PMod}_{g,
b,p+1} \to \on{PMod}_{g,0,b+p+1}$.
The action of $\Gamma'$ on
$H_1(\Sigma_{g',0,b'+p'}, \mathbb{C})$ factors through the natural map $\Gamma'\to \Gamma''$, inducing the claimed bijection.
\end{proof}
%\begin{lemma}\label{lemma:versal-H-covers}
%As in \autoref{notation:prym}, let $\phi: \pi_1(\Sigma_{g, 0, n}, v_0)\twoheadrightarrow H$ be a surjection. There exists the data of
%\begin{enumerate}
%\item a dominant \'etale morphism $\mathscr{M}\to \mathscr{M}_{g,n}$, with $\pi^\circ: \mathscr{C}^\circ\to \mathscr{M}$ the associated family of punctured curves,
%\item a point $c\in \mathscr{C}^\circ$, $m=\pi^\circ(c)$, and an identification $i: \pi_1(\Sigma_{g, 0, n}, v_0)\simeq \pi_1(\mathscr{C}_m^\circ, c)$, and
%\item a finite \'etale Galois $H$-cover $f: \mathscr{X}^\circ\to \mathscr{C}^\circ$
%\end{enumerate}
%such that under the identification $i$, $f^{-1}(\mathscr{C}_m^\circ)\to \mathscr{C}_m^\circ$ is the $H$-cover of $\mathscr{C}_m^\circ$ corresponding to $\phi$.
%\end{lemma}
%We refer to the data produced by \autoref{lemma:versal-H-covers} as a \emph{versal family of $H$-covers}.
%\begin{proof}[Proof of \autoref{lemma:versal-H-covers}]
%%As in \autoref{notation:prym}, we let $\Gamma\subset \on{PMod}_{g, 0, n+1}$ be the subgroup preserving $\phi$ up to conjugacy, and let $\widetilde{\Gamma}$ be its image in $\on{PMod}_{g, 0, n}$. 
%
%Let $\{\rho_i\}$ be the set of irreducible representations of $H$. Combining \autoref{lemma:geometric-mcg-rep-construction} and \autoref{proposition:gl-lift}, we see that for each $i$ there exists a versal family of $n$-punctured curves of genus $g$, $\pi_i^\circ: \mathscr{C}_i^\circ\to \mathscr{M}_i$, and a local system $\mathbb{V}_i$ on $\mathscr{C}_i^\circ$ with fibral monodromy given by $\rho_i$. Replacing the $\mathscr{M}_i$ by a component of their fiber product over $\mathscr{M}_{g,n}$, we may assume all the $\mathscr{M}_i$ are equal, so we have a versal family $\pi^\circ:{\mathscr{C}'}^\circ\to \mathscr{M'}$ carrying local systems $\mathbb{V}_i$.
%
%Let $\mathscr{X}^\circ$ be the cover of $\mathscr{C}^\circ$ defined by the joint kernel of the monodromy representations on the $\mathbb{V}_i$. The map $\mathscr{X}^\circ\to \mathscr{M}'$ may not have connected fibers; let $\mathscr{X}^\circ \to \mathscr{M}\to \mathscr{M}'$ be the Stein factorization, and let $\mathscr{C}^\circ=\mathscr{M}\times_{\mathscr{M}'}{\mathscr{C}'}^\circ$. As the set of irreducible representations of a finite group are jointly faithful, this cover has the desired properties.
%%
%%For each $\gamma$, any two choices of $h_\gamma$ differ by an element of the center $Z(H)$ of $H$. Hence the obstruction to lifting the set-theoretic map $$h: \gamma\mapsto h_\gamma$$ to a group homomorphism $\Gamma\to H$ is an element $o$ of $H^2(\Gamma, Z(H))$. Let $\mathscr{M}_{\widetilde\Gamma}$ be the cover of $\mathscr{M}_{g,n}$ corresponding to $\widetilde\Gamma$; by shrinking $\widetilde\Gamma$, we may assume $\mathscr{M}_\Gamma$ is a scheme. 
%%
%%
%\end{proof}
\subsection{}
	\label{subsection:proof-putman-wieland}
\source{canonical-representations-surface-groups}

We now prove 
\autoref{theorem:asymptotic-putman-wieland}, which
is more or less an immediate consequence of \autoref{theorem:rank-bound-subsystem}
applied to the irreducible representations of $H$,
once one sets up the appropriate local systems.
For the proof, we will need the existence of a versal family of $\phi$-covers,
for $\phi: \pi_1(\Sigma_{g, 0, n}, v_0)\twoheadrightarrow H$,
which we now define.
\begin{definition}
\source{canonical-representations-surface-groups}
\notready
	\label{definition:}
\source{canonical-representations-surface-groups}
	As in \autoref{notation:prym}, specify a surjection
	$\phi: \pi_1(\Sigma_{g, 0, n}, v_0)\twoheadrightarrow H$.
	A {\em versal family of $\phi$-covers} is the data of 
	\begin{enumerate}
\item a dominant \'etale morphism $\mathscr{M}\to \mathscr{M}_{g,n}$, with $\pi^\circ: \mathscr{C}^\circ\to \mathscr{M}$ the associated family of punctured curves,
\item a point $c\in \mathscr{C}^\circ$, $m=\pi^\circ(c)$, $C^\circ = \mathscr
C^\circ_m$, and an identification $i: \pi_1(\Sigma_{g, 0, n}, v_0)\simeq
\pi_1(C^\circ, c)$, and
\item a finite \'etale Galois $H$-cover $f: \mathscr{X}^\circ\to
	\mathscr{C}^\circ$ inducing a map
	$$\pi_1(C^\circ, c) \to \pi_1(\mathscr{C}^\circ,
	c)/\pi_1(\mathscr{X}^\circ,x) \simeq H$$
	agreeing with the surjection $\phi$ under the
	identification of $(2)$.
\end{enumerate}
\end{definition}

\begin{proof}[Proof of \autoref{theorem:asymptotic-putman-wieland}]
	Using \autoref{lemma:ignore-boundary}, 
	it is enough to prove
	\autoref{theorem:asymptotic-putman-wieland}
	when $b =0$. 

	By \cite[Theorem 4]{wewers:thesis}, there exists a versal family of
	$\phi$-covers $\mathscr X^\circ \xrightarrow{f} \mathscr C^\circ
	\xrightarrow{\pi^\circ} \mathscr M$.
	Technically, 
	\cite[Theorem 4]{wewers:thesis} constructs $\mathscr M$ as a Deligne-Mumford stack,
	where points have isotropy groups isomorphic to the center of $H$,	
	but we may replace this stack by
	a scheme which has a dominant \'etale map to this stack. 
	Further, after replacing $\mathscr M$ with another scheme that has a dominant \'etale map
	to $\mathscr M$, corresponding generically to an \'etale multisection of $\pi^\circ$, we may assume $\pi^\circ$ admits a section, and hence that the map $\mathscr{M}\to \mathscr{M}_{g,n}$ lifts to a map $\mathscr{M}\to \mathscr{M}_{g,n+1}$.

	As $f$ is a Galois finite \'etale $H$-cover, for each irreducible
	representation $\rho$ of $H$, we obtain a local system $\mathbb{U}_\rho$
	on $\mathscr{C}^\circ$, with monodromy representation given by $\rho$,
	and with $\mathbb{U}_\rho|_{\mathscr{X}^\circ}$ trivial. After replacing $\mathscr{M}$ with a dominant \'etale cover, we may write
	$$f_*\mathbb{C}=\bigoplus_\rho \mathbb{U}_\rho\otimes
	(\pi^\circ)^*\mathbb{W}_\rho,$$ where
	$\mathbb{W}_\rho:=\pi^\circ_*\on{Hom}(\mathbb{U}_\rho, f_*\mathbb{C})$,
	by \autoref{lemma:unitary-direct-sum} (where we take $A=\mathbb{C}$). 
	As
	$\mathbb{U}_\rho$ has finite monodromy by
	\autoref{lemma:finite-rho-implies-finite-lift}(1), and $f_*\mathbb{C}$
	has finite monodromy by definition, the same is true for
	$\mathbb{W}_\rho$. Thus after replacing $\mathscr{M}$ with a finite \'etale cover we may assume each
	$\mathbb{W}_\rho$ is a trivial local system.
	
	We next verify the first part of
	\autoref{theorem:asymptotic-putman-wieland}, about $H^1(\Sigma_{g', 0,
	p'}, \mathbb{Q})^\rho$.
	Observe that for $m\in \mathscr{M}$, the action of $\pi_1(\mathscr{M},
	m)$ on $$R^1\pi^\circ_*( \mathbb{U}_\rho\otimes
	(\pi^\circ)^*\mathbb{W}_\rho)_m=(R^1\pi^\circ_* \mathbb{U}_\rho\otimes
	\mathbb{W}_\rho)_m$$ is identified with its action on
	$H^1(\mathscr{X}^\circ_m, \mathbb{C})^\rho=H^1(\Sigma_{g', 0, p'},
	\mathbb{C})^\rho$. 
	Note also that the action of
	$\pi_1(\mathscr{M})$ factors through $\Gamma$, 
	under our given map $\pi_1(\mathscr{M})\to \pi_1(\mathscr{M}_{g,n+1})=\on{PMod}_{g,n+1}$.
	To prove the first part of
	\autoref{theorem:asymptotic-putman-wieland}, about $H^1(\Sigma_{g', 0,
	p'}, \mathbb{Q})^\rho$, it thus suffices to show that, after replacing
	$\mathscr{M}$ with an arbitrary finite \'etale cover, $R^1(\pi^\circ_*
	\mathbb{U}_\rho)\otimes \mathbb{W}_\rho$ contains no non-trivial
	sub-local systems of rank less than $2g-2\dim\rho$. 
	As $\mathbb{W}_\rho$
	is a trivial local system, this holds by applying
	\autoref{theorem:rank-bound-subsystem} to $\mathbb{U}_\rho$.

We have  thus far proven a statement for the cohomology of $\Sigma_{g',n'}$, and to
conclude we deduce a corresponding statement for the homology of $\Sigma_{g'}$.
Using the inclusion 
$H^1(\Sigma_{g'}, \mathbb C)^\rho \subset H^1(\Sigma_{g', n'}, \mathbb{C})^\rho$
we deduce that the former can have no nonzero $\Gamma'$-invariant subrepresentations 
(for $\Gamma' \subset \Gamma$ finite index)
of dimension less than $2g - 2\dim \rho$, as the latter has no such 
subrepresentations.
Using Poincar\'e duality and the intersection pairing on $H_1(\Sigma_{g'}, \mathbb C)$ 
we obtain a $\on{Mod}_{g,n}$ equivariant isomorphism
$H^1(\Sigma_{g'},\mathbb C) \simeq H_1(\Sigma_{g'},\mathbb C)$.
Hence, $H_1(\Sigma_{g'}, \mathbb C)^\rho$ also has no subrepresentations of 
dimension less than $2g - 2\dim\rho$.
\end{proof}
